\documentclass[../../../../main.tex]{subfiles}

\begin{document}

\section*{第1問}

{\bf [解]}
$L:y=2x\ (0\le x\le1)$である．よって，
\begin{align*}
x^2+ax+b=2x\quad(0\le x\le1)
\end{align*}
が解を持てば良い．$f(x)=x^2+(a-2)x+b$とおく．

\underline{$1^\circ$ 区間内に$1$つ}
\begin{align*}
f(0)f(1)\le0\iff b(a+b-1)\le0
\end{align*}

\underline{$2^\circ$ 区間内に$2$つ（重解含む）}
\begin{align*}
\text{判別式：}&(a-2)^2-4b\ge0\\
\text{端点：}&f(0)\ge0，\ f(1)\ge0\\
\text{軸：}&0\le\frac{2-a}2\le1
\end{align*}
\begin{align*}
\iff\quad b\le\frac14(a-2)^2，\qquad b\ge0，\ a+b-1\ge0，\qquad0\le a\le2
\end{align*}

以上を図示して，下図斜線部（境界含む）．

\begin{figure}[htb]
\centering
\begin{tikzpicture}[scale=1.1, >=stealth]
  \begin{axis}[
      axis lines=middle,
      xlabel={$a$}, ylabel={$b$},
      xmin=-2.5, xmax=3.5, ymin=-1.5, ymax=3.2,
      domain=-2.5:3.5, samples=100, smooth,
      clip=false,
    ]
    \addplot[name path=parab, domain=0:2, thick] {0.25*(x-2)^2};
    \addplot[name path=lineb, thick] {1-x};
    \addplot[name path=zero, domain=-2.5:3.5] {0};
    \addplot[gray!25] fill between[of=lineb and zero, soft clip={domain=-2.5:1}];
    \addplot[gray!25] fill between[of=parab and lineb, soft clip={domain=0:2}];
    \node[right] at (axis cs:2,1.3) {$b=\frac14(a-2)^2$};
    \node[below right] at (axis cs:2.2,-1.2) {$b=1-a$};
  \end{axis}
\end{tikzpicture}
\caption{条件を満たす$(a,b)$の範囲（斜線部，境界含む）}
\end{figure}

\newpage

\section*{第2問}

{\bf [解]}
各辺正から常用対数をとって，$\log_{10}x=\log x$とおくと，
\begin{align*}
10\log2<n\log\frac54<20\log2
\end{align*}
\begin{align*}
10\log2<n\cdot\log\frac{10}8<20\log2
\end{align*}
\begin{align*}
10\log2<n(1-3\log2)<20\log2
\end{align*}

$\alpha=\log2$とおく．$1-3\alpha>0$から，
\begin{align*}
\frac{10\alpha}{1-3\alpha}<n<\frac{20\alpha}{1-3\alpha}
\end{align*}
\begin{align*}
-\frac{10}3+\frac{10}{3(1-3\alpha)}<n<\frac{20}3\left(-1+\frac1{1-3\alpha}\right)\quad\cdots①
\end{align*}

ここで，$0.301<\alpha<0.3011$から，
\begin{align*}
0.0967<1-3\alpha<0.097
\end{align*}
\begin{align*}
\frac1{0.097}<\frac1{1-3\alpha}<\frac1{0.0967}
\end{align*}
\begin{align*}
10.30\ldots<\frac1{1-3\alpha}<10.34\ldots
\end{align*}
\begin{align*}
9.30\ldots<-1+\frac1{1-3\alpha}<9.34\ldots
\end{align*}

となるので，
\begin{align*}
31.0\ldots<\frac{10}3\left(-1+\frac1{1-3\alpha}\right)<31.13\ldots
\end{align*}
\begin{align*}
62.0\ldots<\frac{20}3\left(-1+\frac1{1-3\alpha}\right)<62.26\ldots
\end{align*}

だから，①をみたす$n$は，$n=32,\ldots,62$の$62-32+1=31$個．

\newpage

\section*{第3問}

{\bf [解]}
\begin{align*}
\beta\gamma+\gamma\alpha+\alpha\beta=\frac12\left[(\alpha+\beta+\gamma)^2-(\alpha^2+\beta^2+\gamma^2)\right]=0
\end{align*}

だから，$\alpha\beta\gamma=k^3$とおくと，$\alpha,\beta,\gamma$は，$x$の3次方程式$x^3-k^3=0$の3解で，$\alpha\ne\beta$，$\gamma\ne\alpha$かつ$\alpha\ne0$である．
\begin{align*}
(x-k)(x^2+kx+k^2)=0
\end{align*}
\begin{align*}
\therefore\ x=k,\ \frac{-1\pm\sqrt3i}2k
\end{align*}

だから，$\triangle ABC$は正三角形（重心は原点）．

\newpage

\section*{第4問}

{\bf [解]}
$217=7\cdot31$である．
\begin{align*}
(a-b)(a^2+ab+b^2)=7\cdot31\quad\cdots①
\end{align*}

$a^2+ab+b^2>0$から，$a-b>0$で，
\begin{align*}
(a-b,\ a^2+ab+b^2)=(1,217),\ (7,31),\ (31,7),\ (217,1)
\end{align*}

のうち，整数解を持つのは最初の2つのみで（後の2つは判別式が負），
\begin{align*}
(a,b)=(9,8),\ (-8,-9),\ (6,-1),\ (1,-6)．
\end{align*}

\newpage

\section*{第5問}

{\bf [解]}
$k\in\mathbb{Z}_{\ge1}$とする．

\textbf{(1)} $y=\cos x$は区間内で単調減少で$-1\le\cos x\le1$である．一方，
\begin{align*}
0\le\frac1{x^2+1}\le1\iff-1\le-1+\frac2{x^2+1}\le1\iff-1\le\frac{1-x^2}{1+x^2}\le1
\end{align*}
である．したがって，$f(x)=\cos x-\dfrac{1-x^2}{1+x^2}$とすると，
\begin{align*}
f(2k\pi)\ge0，\qquad f((2k+1)\pi)\le0
\end{align*}
となるから，$f$が連続であることも用いて，$f(\alpha_k)=0\ (2k\pi\le\alpha_k\le(2k+1)\pi)$なる$\alpha_k$がある（$C_1$，$C_2$の共有点の存在）．$C_1$の接線$\ell_k$は
\begin{align*}
\ell_k:y=-\sin\alpha_k(x-\alpha_k)+\cos\alpha_k=g(x)
\end{align*}
で表せる．$f(\alpha_k)=0$から
\begin{align*}
\cos\alpha_k=\frac{1-\alpha_k^2}{\alpha_k^2+1}，\qquad\sin\alpha_k=\sqrt{1-\frac{(1-\alpha_k^2)^2}{(\alpha_k^2+1)^2}}=\frac{2\alpha_k}{\alpha_k^2+1}\ (>0)
\end{align*}
であり，
\begin{align*}
g(0)=-\frac{2\alpha_k}{\alpha_k^2+1}(-\alpha_k)+\frac{1-\alpha_k^2}{\alpha_k^2+1}=\frac{\alpha_k^2+1}{\alpha_k^2+1}=1
\end{align*}
から，$\ell_k$は$(0,1)$を通る．$\blacksquare$

\textbf{(2)} $h(x)=g(x)-\dfrac{1-x^2}{1+x^2}$とおく．
\begin{align*}
h(x)&=-\frac{2\alpha_k}{\alpha_k^2+1}x+1-\frac{1-x^2}{1+x^2}\\
&=-\frac{2\alpha_k}{\alpha_k^2+1}x+2-\frac2{1+x^2}\\
&=\frac{-2x}{(\alpha_k^2+1)(x^2+1)}(x-\alpha_k)(\alpha_kx-1)
\end{align*}

$2k\pi\le x,\alpha_k\le(2k+1)\pi$（$k\ge1$）では$x>0$，$\alpha_kx-1>0$だから，$h(x)$の符号は$-(x-\alpha_k)$の符号と一致し，
\begin{align*}
x<\alpha_k\text{の時，}h(x)>0，\qquad x=\alpha_k\text{で，}h(x)=0，\qquad\alpha_k<x\text{の時，}h(x)<0
\end{align*}
となり，たしかに$2k\pi\le x\le(2k+1)\pi$では$h(x)=0$はただ1つの解を持つ．つまり，$C_1$と$C_2$の共有点はただ1つである．$\blacksquare$

\bigskip
\noindent{\bf [解2]}
\textbf{(2)} $x=t$における$C_1$の接線は$y=-\sin t(x-t)+\cos t$であり，これより，切片の関数$Y(t)=t\sin t+\cos t$，$Y'(t)=t\cos t$から下表を得る．

\begin{center}
\begin{tabular}{c|ccccc}
$t$ & $2k\pi$ & & $\left(2k+\frac12\right)\pi$ & & $(2k+1)\pi$ \\
\hline
$Y'$ & & $+$ & $0$ & $-$ & \\
\hline
$Y$ & $1$ & $\nearrow$ & & $\searrow$ & $-1$ \\
\end{tabular}
\end{center}

$\alpha_k\ne2k\pi$であることと，$2k\pi$から$(2k+1)\pi$で$Y(t)=1$なる$t$がただ1つ（$\alpha_k$自身）であることから，これと(1)をあわせて，題意は示された．$\blacksquare$

\newpage

\section*{第6問}

{\bf [解]}
左から$k$両目まで見た時，$k$両目が赤である場合の数を$a_k$，青か黄色である場合の数を$b_k$とする．題意から
\begin{align*}
a_{k+1}&=a_k+b_k\\
b_{k+1}&=2a_k
\end{align*}
であり，又，$a_1=1$，$b_1=2$となる．※から$b_k$を消して，
\begin{align*}
a_{k+2}=a_{k+1}+2a_k，\qquad a_1=1，\ a_2=3
\end{align*}

だから，解いて，
\begin{align*}
a_k&=\frac13\left\{2^{k+1}+(-1)^k\right\}\\
b_k&=\frac23\left\{2^k+(-1)^{k-1}\right\}\quad(k\ge2)
\end{align*}

だから$n\ge2$の時，もとめるのは
\begin{align*}
a_n+b_n&=\frac13\left(2^{n+2}+(-1)^n+2(-1)^{n-1}\right)\\
&=\frac13\left(2^{n+2}+(-1)^{n-1}\right)．
\end{align*}

\newpage

\end{document}
